\documentclass[12pt]{article}

% ----- Packages -----
\usepackage[utf8]{inputenc}
\usepackage[T1]{fontenc}
\usepackage{mathptmx}               % Times New Roman clone
\usepackage[margin=1in]{geometry}
\usepackage{setspace}
\onehalfspacing
\usepackage{hyperref}
\hypersetup{colorlinks=true, linkcolor=blue, citecolor=blue, urlcolor=blue}
\usepackage{natbib}
\bibliographystyle{apalike}
\usepackage{enumitem}
\usepackage{parskip}

% ----- Title -----
\title{\textbf{Project Proposal} \\[6pt]
       \large Case Studies in Machine Learning --- Spring 2026}
\author{Juan P. Madrigal-Cianci, Ph.D \\ UT EID: jm225572}
\date{February 17, 2026}

\begin{document}
\maketitle

% ==================================================================
\section*{Forecasting Prediction Market Prices Using Ensemble and Deep Learning Models}
% ==================================================================

% ------------------------------------------------------------------
\subsection*{1.\quad Introduction and Motivation}
% ------------------------------------------------------------------

Prediction markets, that is, platforms such as Kalshi and Polymarket where participants trade contracts tied to the outcomes of real-world events, have emerged as powerful aggregators of collective belief.
Recent empirical evidence suggests that prediction market prices frequently rival or surpass expert forecasts and traditional polling in domains ranging from elections to macroeconomic indicators \citep{Arrow2008, Wolfers2004}.
Yet despite their informational richness, comparatively little work has applied modern machine learning to the task of \textit{forecasting the price trajectories} of prediction market contracts themselves.
This gap is surprising given the growing volume of publicly available order-book and pricing data, the high-frequency nature of these markets, and the natural suitability of sequence-based deep learning architectures to such time-series problems.

A recent large-scale study by \citet{Clinton2025} analyzed over 2{,}500 political contracts across Kalshi, Polymarket, and PredictIt during the 2024 U.S.\ presidential election, revealing significant cross-platform price disparities and efficiency failures; suggesting that ML-based forecasting models could exploit predictable patterns in contract price dynamics.
Similarly, \citet{ng2025price} provide the first evidence on price discovery dynamics across modern prediction markets, finding that net order imbalance from large trades strongly predicts subsequent returns.

This project proposes a systematic comparison of classical supervised learning models and modern deep learning architectures for short-horizon price prediction on prediction market contracts.
The goal is twofold: 
\begin{enumerate}
\item to evaluate which ML methods best capture the dynamics of belief-driven markets, and 
\item to investigate whether external features, such as news sentiment, polling data, and social-media activity, provide statistically significant improvements over price-only baselines.
\end{enumerate}

The findings would contribute to the growing literature at the intersection of financial machine learning and information aggregation, with the potential for publication in venues such as \textit{Finance Research Letters}.

% ------------------------------------------------------------------
\subsection*{2.\quad Research Questions}
% ------------------------------------------------------------------

\begin{itemize}[leftmargin=2em]
  \item \textbf{RQ1:} How accurately can ensemble methods (Random Forest, XGBoost, LightGBM) and recurrent deep learning models (LSTM, GRU, Transformer-based architectures) forecast short-term price movements of prediction market contracts?
  \item \textbf{RQ2:} Does the incorporation of exogenous features (e.g., news sentiment scores, polling aggregates, Google Trends, and social-media volume) improve forecast accuracy relative to price-only (endogenous) baselines?
  \item \textbf{RQ3:} What market microstructure features (e.g., bid--ask spread, volume spikes, time-to-expiration) are most predictive of contract price changes, and how do feature importances differ across model families?
\end{itemize}

% ------------------------------------------------------------------
\subsection*{3.\quad Data Sources}
% ------------------------------------------------------------------

The primary dataset will consist of historical contract-level pricing data from \textbf{Polymarket} (sourced via its public API and the Polygonscan blockchain explorer) and \textbf{Kalshi} (via the Kalshi public market data API).
Contracts will span multiple event categories---political (e.g., U.S.\ elections, policy decisions), economic (e.g., Fed rate decisions, CPI prints), and cultural/sports---to ensure generalizability.
The target sample will cover contracts active between January 2023 and December 2025, yielding thousands of unique contract time series.
Exogenous features will be drawn from news APIs (e.g., NewsAPI, GDELT), polling aggregators (e.g., FiveThirtyEight, RealClearPolitics), Google Trends, and Twitter/X volume data.

% ------------------------------------------------------------------
\subsection*{4.\quad Proposed Methods}
% ------------------------------------------------------------------

The modeling pipeline will follow the structure introduced in class.
After data collection and cleaning, features will be engineered at multiple time scales (hourly, daily) and standardized.
The study will benchmark the following model families:

\begin{itemize}[leftmargin=2em]
  \item \textbf{Baseline models:} Linear Regression, Ridge/Lasso, na\"ive persistence forecast.
  \item \textbf{Tree-based ensembles:} Random Forest, Gradient Boosted Trees (XGBoost, LightGBM), with hyperparameter tuning via Optuna \citep{Chen2016}.
  \item \textbf{Deep learning:} LSTM \citep{Hochreiter1997}, GRU, and a lightweight Temporal Fusion Transformer (TFT) \citep{Lim2021}, implemented in PyTorch.
  \item \textbf{Evaluation:} Models will be evaluated using MAE, RMSE, directional accuracy, and a simulated trading P\&L metric. Walk-forward (expanding window) cross-validation will be used to respect temporal ordering and avoid lookahead bias.
\end{itemize}


\bibliography{bib}


\end{document}
